\documentclass[11pt]{article}
\usepackage{amsmath,amssymb,amsthm}
\newtheorem{theorem}{Theorem}
\newtheorem{proposition}{Proposition}
\begin{document}
\title{A009: fixed-posterior fractional incidence certificates}
\author{A009, stream B02}\maketitle
All entropy units are bits except the explicitly natural-unit program.
This is provisional continuation evidence for A006/A007, pending independent
review. Full FEI is open. Complete source locators and the full ledger are
in \texttt{REPORT.md}.

Fix actual partition-forest laws for a finite nonconstant gate union.
Let $v$ be nonempty spectral mass, $c_l$ expected use of gate $l$,
$\pi_l>0$ its root parameter, and $A_{jl}=1_{j\in R_l}$.
Remove unused gates if necessary. Every retained support is nonempty.
For a constant-false gate removal is mandatory; a constant-true gate makes
the Boolean function constant. Duplicate supports retain separate columns.
The entire weight-dependent residual is
\[
 R(\lambda)=d+\frac{F(\ln\lambda)}{\ln2},\quad
 F(t)=v\ln D(e^t)-\sum_lc_lt_l,\quad
 D(\lambda)=\prod_l(1+\pi_l\lambda_l)-1.
\]
Here $d$ retains all ranks, subset congestion, cover multiplicity, root
rarity, and minus the entropy of the laws on distinct actual runs.

\begin{theorem}[Fixed-law convex primal and dual certificate]
The program minimizing $F(t)$ subject to
$\sum_l A_{jl}e^{t_l}\le1$ is convex. For a probability distribution $\nu$
on nonempty root sets, put $r_l=\Pr_\nu(l\in R)$. For $y_j\ge0$, put
$q_l=\sum_j A_{jl}y_j$ and $b_l=c_l-vr_l$. If $b_l\ge0$, and
$q_l>0$ whenever $b_l>0$, a lower bound on the infimum is
\[
 L(\nu,y)=v\left(H_{\rm nat}(\nu)+\sum_lr_l\ln\pi_l\right)
 -\sum_jy_j+\sum_{l:b_l>0}b_l\left(1+\ln\frac{q_l}{b_l}\right).
\]
For positive feasible $\lambda$, take
\[
 \nu_*(R)=\frac{\prod_{l\in R}\pi_l\lambda_l}{D(\lambda)},\qquad
 r_l=\frac{\pi_l\lambda_l\prod_{k\ne l}(1+\pi_k\lambda_k)}{D(\lambda)}.
\]
If $c_l-vr_l=\lambda_lq_l$ and
$y_j(\sum_l A_{jl}\lambda_l-1)=0$, then $L(\nu_*,y)=F(\ln\lambda)$.
Thus these identities certify exact global optimality.
\end{theorem}
\begin{proof}
Expansion expresses $\ln D(e^t)$ as a log-sum-exp over nonempty subsets,
so $F$ and the constraints are convex. Gibbs gives
$\ln D(e^t)\ge H_{\rm nat}(\nu)+\sum_l r_l(t_l+\ln\pi_l)$.
Add the nonpositive-on-feasible-points quantity
$\sum_jy_j(\sum_l A_{jl}e^{t_l}-1)$ and take the unconstrained infimum.
Each term is $-b_lt_l+q_le^{t_l}$; its infimum is
$b_l(1+\ln(q_l/b_l))$ if $b_l,q_l>0$, and zero if $b_l=0$.
The stated equalities give equality in Gibbs, every scalar infimum, and
the coordinate-price inequality. No minimax interchange is needed.
\end{proof}

For the actual Boolean instance below with $(r,m)=(1,2)$,
$v=15/16$, $c_l=17/32$, $\pi_l=1/4$, $\lambda_l=1/2$,
$r_l=9/17$, and the sole nonzero coordinate price $y=19/272$ satisfy
the theorem by rational arithmetic. Other illustrative abstract
certificates are not claimed Boolean-realizable.

\begin{proposition}[The optimized posterior changes convexity]
For the duplicate-dictator representation $x_1=x_1\vee x_1$, the
posterior-optimized natural-unit residual is
\[
 \ln\left(1+\frac{\lambda_1\lambda_2}{2(\lambda_1+\lambda_2)}\right).
\]
With $\lambda_1=e^{t+s}$, $\lambda_2=e^{t-s}$ this is
$\ln(1+e^t/(4\cosh s))$. Its second $s$ derivative at
$t=\ln(1/4)$, $s=0$ is $-1/17$, in the strict feasible interior.
The joint posterior-optimized objective is therefore nonconvex.
\end{proposition}
\begin{proof}
There are precisely two runs for the singleton spectral set, with
$2^{-B_l}=\lambda_l/2$. The finite Gibbs optimum is minus the
logarithm of their sum. Divide the expanded root normalizer $D$
by this sum and differentiate the displayed formula.
\end{proof}

\begin{theorem}[Complete shared-core singleton certificate]
For every $r,m\ge1$, let $f=2A_rB_m-1$, where $A_r$ is a core
conjunction and $B_m$ an OR of private literals on disjoint fair inputs.
Use the $m$ distinct conjunction gates $A_r\wedge x_l$, with core-first
orders, one cover per gate, and $\lambda_l=1/m$. The conditioned
posterior forest residual admits actual laws with
\[
 R\le 8 I(f).
\]
A normalized prior on complete scaled references can be paid with
$R+\text{prior cost}\le10I(f)$. These functions already have read-once
formula representations; the result certifies their DNF of read $m$.
\end{theorem}
\begin{proof}
Write $p=2^{-r}$, $q=2^{-m}$, $b=1-q$, $\gamma=6/5$.
The exact law on a nonempty set $U\cup T$ is $4p^2b^2$ if
$T=\varnothing$, and $4p^2q^2$ otherwise. Consequently
\[
 v=4pb(1-pb),\quad I_c=2rpb,\quad I_d=2pmq.
\]
For nonempty $U$ and empty $T$, average over all $m$ single roots.
For both nonempty, average over the $|T|!$ actual star runs, choosing
the root among touched private gates and all child-label orders.
For empty $U$, all touched private gates are forced roots and the run
is unique. Every selected subset is disjoint, and every used gate is
visited once. The star root reaches all used gates through a selected
core coordinate, so these are valid source-component forests.

The root count $K$ obeys
$E K-v=4p^2q(m/2-b)\le pI_d$.
Since $D=(1+p/(2m))^m-1$, expansion and the exponential integral give
$\log_2D=-(r+1)+\delta$ with $0\le\delta\le(p/2)\log_2 e$.
Thus the normalization and root rarity together cost at most
$I_d+(\log_2e)I_c$, using $(r+1)p\le1$ and $v\le4pb$.

For a selected core rank $j$, its predecessor count is
$\operatorname{Bin}(j-1,1/2)$, and
$E[1/(1+\operatorname{Bin})]=2(1-2^{-j})/j$.
Jensen implies expected log-rank discrepancy at most one bit
conditional on selection. Summing its contribution under the exact
Fourier law pays all core ranks by $\gamma I_c$.
For each private leaf its congestion and rank terms combine into
$\gamma\log_2(r+1)$; the inside root has no larger cost.
They sum to at most $\gamma\log_2(r+1)I_d\le\gamma I_c$, since
$m2^{-m}\le1/2$, $b\ge1/2$, and $\log_2(r+1)\le r$.
The core-only posterior entropy cancels the label cost exactly.
The remaining labels minus entropy cost at most
$I_d\log_2m\le4p\le4I_c$, by $\log_2m\le m$ and $m^2/2^m\le2$.
All terms have been retained, so
\[
 R\le(4+2\gamma+\log_2e)I_c+I_d<8(I_c+I_d).
\]
Finally weights $w_j=1/[(j+1)(j+2)]$ on complete references with
$\lambda_l^{(j)}=2^{-j}/m$ sum to one. Choosing $j=0$ pays $v\le2I$.
\end{proof}
\section*{Lead-directed outer-pressure test}
The B02 lead assigned the outer fractional pressure
$P=(48/5)\sum_j\alpha_j\log_2 S_j$, where
$S_j=\sum_l\sqrt{q_{j,l}}\lambda_l^{-11/48}$.
For the one-core sunflower, the first-positive-index tie gives
$q_{0,l}=2^{-l}/(1-2^{-m})$, and each private derivative has a
unique gate. The global vector $\lambda_l=2^{-l}$ obeys all rows.
Its common sum is at most $\sqrt2/(2^{13/48}-1)<9$, and its private
log costs are $11l/48$. Thus $P<41I$, a calibration on a known FEI class.

\begin{proposition}[Affine-plane pressure obstruction]
For every prime power $Q\ge16$, let $f_Q$ be the OR of all line
conjunctions in the affine plane over $\mathbb F_Q$. These are its
prime implicants. For any disjoint pivotal tie decomposition and every
global feasible incidence vector,
\[
 P\ge7 I(f_Q)\log_2\frac{Q+1}{2}.
\]
Uniform weights give $P\le7I(f_Q)\log_2(Q+1)$.
\end{proposition}
\begin{proof}
There are $Q(Q+1)$ lines and $Q+1$ through each point. Fix a point $j$
and condition on the other $Q-1$ points of one incident line being
positive. Each alternative incident line is ready with probability
$2^{-(Q-1)}$; each nonincident line has at most one forced point,
so its full-positive probability is at most $2^{-(Q-1)}$.
The union bound leaves clean probability at least $1-\epsilon_Q$,
where $\epsilon_Q=Q(Q+1)2^{-(Q-1)}<1/2$.
Clean events for distinct incident lines are disjoint and make $j$
pivotal. Thus $\alpha_j\ge(Q+1)2^{-(Q-1)}(1-\epsilon_Q)$.
Any tie piece is supported where its incident line is ready, so its
normalized squared mass is at most $2/(Q+1)$.
Write $k=Q+1$, $a=11/48$, $\beta=24/59$.
H\"older and $\sum_l\lambda_l\le1$ give
\[
 S_j\ge\left(\sum_l q_{j,l}^{\beta}\right)^{1+a}
 \ge(k/2)^{35/48}.
\]
Use $q_l^\beta=(\sqrt{q_l}\lambda_l^{-a})^{1/(1+a)}
\lambda_l^{a/(1+a)}$ for the first step, and
$q_l^\beta\ge q_l(2/k)^{\beta-1}$ for the second.
Multiply by $(48/5)\alpha_j$ and sum. For uniform $\lambda=1/k$,
Cauchy gives $\sum\sqrt q\le\sqrt k$, proving the upper bound.
The rational bound on $\epsilon_Q$ holds at $Q=16$ and decreases
with ratio $(Q+2)/(2Q)<1$ thereafter.
\end{proof}
This is a lower bound for the positive coordinatewise pressure only,
not for the full expression with $v\log Z$, whole-set forest cost,
PI-J, or actual Fourier entropy divided by influence.

The finite Walsh and rational-log checks are supporting certificates,
not the all-parameter proof. Arbitrary overlapping base gates do not
have the binomial/private-label law just used. Paying their full
expected structural cost remains the universal bridge.
\end{document}
